\documentclass[10pt,a4paper]{amsart}
\usepackage{amsmath,amssymb,amsthm,graphicx}
% Retained minimal preamble for the Stacks Project chapter more-algebra.tex (More on Algebra).
% The chapter file pinned-original-source.tex carries only \input{preamble} on its first line, so no
% preamble text lives inside the chapter itself.  The Stacks Project uses one single shared preamble for
% every chapter; the snapshot stages it identically as sources/stacks-chow/preamble.tex and
% sources/stacks-etale-cohomology/preamble.tex.  The block below is copied mechanically, line for line,
% from sources/stacks-chow/preamble.tex lines 14-18 (the Xy-pic setup) and 154-217 (the native theorem
% declarations together with the mathematical macros).  Nothing was added, removed or redefined; the
% wrapper class and the amsmath/amssymb/amsthm/graphicx base packages are supplied by the assembler.
% Chapter: more-algebra.tex -- More on Algebra.
% For commutative diagrams we use Xy-pic
\usepackage[all]{xy}

% We use 2cell for 2-commutative diagrams.
\xyoption{2cell}

% Theorem environments.
%
\theoremstyle{plain}
\newtheorem{theorem}[subsection]{Theorem}
\newtheorem{proposition}[subsection]{Proposition}
\newtheorem{lemma}[subsection]{Lemma}

\theoremstyle{definition}
\newtheorem{definition}[subsection]{Definition}
\newtheorem{example}[subsection]{Example}
\newtheorem{exercise}[subsection]{Exercise}
\newtheorem{situation}[subsection]{Situation}

\theoremstyle{remark}
\newtheorem{remark}[subsection]{Remark}
\newtheorem{remarks}[subsection]{Remarks}

\numberwithin{equation}{subsection}

% Macros
%
\def\lim{\mathop{\mathrm{lim}}\nolimits}
\def\colim{\mathop{\mathrm{colim}}\nolimits}
\def\Spec{\mathop{\mathrm{Spec}}}
\def\Hom{\mathop{\mathrm{Hom}}\nolimits}
\def\Ext{\mathop{\mathrm{Ext}}\nolimits}
\def\SheafHom{\mathop{\mathcal{H}\!\mathit{om}}\nolimits}
\def\SheafExt{\mathop{\mathcal{E}\!\mathit{xt}}\nolimits}
\def\Sch{\mathit{Sch}}
\def\Mor{\mathop{\mathrm{Mor}}\nolimits}
\def\Ob{\mathop{\mathrm{Ob}}\nolimits}
\def\Sh{\mathop{\mathit{Sh}}\nolimits}
\def\NL{\mathop{N\!L}\nolimits}
\def\CH{\mathop{\mathrm{CH}}\nolimits}
\def\proetale{{pro\text{-}\acute{e}tale}}
\def\etale{{\acute{e}tale}}
\def\QCoh{\mathit{QCoh}}
\def\Ker{\mathop{\mathrm{Ker}}}
\def\Im{\mathop{\mathrm{Im}}}
\def\Coker{\mathop{\mathrm{Coker}}}
\def\Coim{\mathop{\mathrm{Coim}}}

% Boxtimes
%
\DeclareMathSymbol{\boxtimes}{\mathbin}{AMSa}{"02}

%
% Macros for moduli stacks/spaces
%
\def\QCohstack{\mathcal{QC}\!\mathit{oh}}
\def\Cohstack{\mathcal{C}\!\mathit{oh}}
\def\Spacesstack{\mathcal{S}\!\mathit{paces}}
\def\Quotfunctor{\mathrm{Quot}}
\def\Hilbfunctor{\mathrm{Hilb}}
\def\Curvesstack{\mathcal{C}\!\mathit{urves}}
\def\Polarizedstack{\mathcal{P}\!\mathit{olarized}}
\def\Complexesstack{\mathcal{C}\!\mathit{omplexes}}
% \Pic is the operator that assigns to X its picard group, usage \Pic(X)
% \Picardstack_{X/B} denotes the Picard stack of X over B
% \Picardfunctor_{X/B} denotes the Picard functor of X over B
\def\Pic{\mathop{\mathrm{Pic}}\nolimits}
\def\Picardstack{\mathcal{P}\!\mathit{ic}}
\def\Picardfunctor{\mathrm{Pic}}
\def\Deformationcategory{\mathcal{D}\!\mathit{ef}}

% --- appended 2026-10-12 (Chief): the project own definitions of the reference/slogan/history
% display environments, copied mechanically from the shared Stacks preamble lines 8-12
% (\usepackage{verbatim} plus the three \newenvironment lines, which the chapter preamble keeps
% verbatim from the project preamble). Nothing added, removed or redefined: these environments
% carry no mathematical content and the project itself typesets them as comments. Without these
% three lines a counted statement whose own text uses them fails with Environment slogan undefined.
\usepackage{verbatim}
\newenvironment{reference}{\comment}{\endcomment}
\newenvironment{slogan}{\comment}{\endcomment}
\newenvironment{history}{\comment}{\endcomment}
\usepackage{cleveref}
\crefname{originallemma}{original lemma}{original lemmas}
\Crefname{originallemma}{Original lemma}{Original lemmas}
\makeatletter
\newlabel{derived-remark-truncation-distinguished-triangle}{{08J5}{source}{}{sourcepointer-1}{}}
\newlabel{derived-remark-truncation-distinguished-triangle@cref}{{[originallemma][0][]08J5}{[1][0][]source}{}{}{}}
\makeatother
\crefname{originallemma}{original lemma}{original lemmas}
\Crefname{originallemma}{Original lemma}{Original lemmas}
\makeatletter
\newlabel{lemma-check-tor-dimension-modulo-nilpotent}{{17307}{source}{}{sourcepointer-2}{}}
\newlabel{lemma-check-tor-dimension-modulo-nilpotent@cref}{{[originallemma][0][]17307}{[1][0][]source}{}{}{}}
\makeatother
\title{Vanishing of the positive cohomology of the reduction modulo an ideal spreads to all power-torsion and bounded complexes with power-torsion cohomology}
\author{The Stacks Project Authors}
\date{stacks-project:more-algebra.tex:local-snapshot-20261012; GNU FDL 1.2 or later}
\begin{document}
\maketitle
\setcounter{subsection}{1}
\noindent\textit{Editorial note.} All delivered mathematical text is the author's own, copied from the pinned Stacks Project chapter file more-algebra.tex; no statement or proof step was rewritten, reordered or removed.\par\medskip
\noindent\textit{Editorial note.} The chapter file pinned-original-source.tex carries only \textbackslash{}input\{preamble\} on its first line, so no preamble text lives inside the chapter. The retained file sources/stacks-more-algebra/retained-preamble.tex is a mechanical, line-for-line copy of the project's single shared preamble as staged in the snapshot; nothing was added, removed or redefined and no package was added.\par\medskip
\noindent\textit{Editorial note.} The source URL is the source-level Stacks Project URL, as required when more than one card is produced from one source; the per-claim tag is carried by claim\_key.\par\medskip
\noindent\textit{Editorial note.} The author writes '(details omitted; hint: represent K by a K-flat complex and compute directly)' and, at the end, 'Details omitted.'; both are kept verbatim. The two references inside the counted proof span are recorded as non-load-bearing navigation pointers.\par\medskip
\noindent\textit{Allowed standard prerequisites (editorial).}\begin{itemize}
\item Derived tensor product as a homotopy colimit over the union of the submodules N[I\textasciicircum{}n]. used at proof lines 24125-24128. Reduces the arbitrary power-torsion module to one annihilated by a fixed power of I.
\item Tor amplitude of a complex and its behaviour modulo a nilpotent ideal. invoked at proof lines 24142-24144. Is the tool that transfers the vanishing from R/I to the square-zero extension.
\item Truncation distinguished triangles and induction on the number of nonzero cohomology modules. used at proof lines 24148-24151. Carries the result from modules to bounded complexes.
\end{itemize}\medskip
\ifdefined\hypertarget\hypertarget{sourcepointer-1}{}\fi
\noindent\textit{Source reference.} Original lemma 08J5: focused-seeds/AGP/sources/stacks-snapshot-20261011/tags.txt, source lines 10838--10838. Cross-chapter navigation pointer into the Stacks Project; the standard fact it supplies is applied inside the counted proof span and the argument resting on it stays there.\par\medskip
\ifdefined\hypertarget\hypertarget{sourcepointer-2}{}\fi
\noindent\textit{Source reference.} Original lemma 17307: sources/stacks-more-algebra/pinned-original-source.tex, source lines 17307--17307. Same-chapter navigation pointer; the statement it names is standard for the section and the argument that consumes it stays inside the counted proof span.\par\medskip
\begin{lemma}
\label{lemma-bounded-above-mod-I}
Let $R$ be a ring. Let $I \subset R$ be an ideal. Let $K$ be an object of
$D(R)$ such that $H^i(K \otimes_R^\mathbf{L} R/I) = 0$ for $i > 0$. Then
\begin{enumerate}
\item $H^i(K \otimes_R^\mathbf{L} R/I^n) = 0$ for all $n \geq 1$
and $i > 0$,
\item $H^i(K \otimes_R^\mathbf{L} N) = 0$ for any $I$-power torsion
$R$-module $N$ and $i > 0$, and
\item for any $M \in D^b(R)$ whose cohomology modules $H^i(M)$
are $I$-power torsion and $0$ for $i > 0$ we have
$H^i(K \otimes_R^\mathbf{L} M) = 0$ for $i > 0$.
\end{enumerate}
\end{lemma}

\begin{proof}
Proof of (2). We can write $N = \bigcup N[I^n]$. We have
$K \otimes_R^\mathbf{L} N = \text{hocolim}_n K \otimes_R^\mathbf{L} N[I^n]$
as tensor products commute with colimits (details omitted; hint: represent
$K$ by a K-flat complex and compute directly).
Hence we may assume $N$ is annihilated by $I^n$.
Consider the $R$-algebra $R' = R/I^n \oplus N$
where $N$ is an ideal of square zero. It suffices to show that the object
$K' = K \otimes_R^\mathbf{L} R'$ of $D(R')$ has vanishing cohomology
in positive degrees.
We have a surjection $R' \to R/I$ of $R$-algebras whose kernel $J$
is nilpotent (any product of $n$ elements in the kernel is zero).
We have
$$
K \otimes_R^\mathbf{L} R/I =
(K \otimes_R^\mathbf{L} R') \otimes_{R'}^\mathbf{L} R/I =
K' \otimes_{R'}^\mathbf{L} R/I
$$
By assumption the complex $K \otimes_R^\mathbf{L} R/I$ has tor-amplitude
in $[-\infty, 0]$. Thus the conclusion by
Lemma \ref{lemma-check-tor-dimension-modulo-nilpotent}.

\medskip\noindent
Part (1) follows trivially from part (2).
Part (3) follows from part (2), induction on the number of nonzero
cohomology modules of $M$, and the distinguished triangles of truncation
from Derived Categories, Remark
\ref{derived-remark-truncation-distinguished-triangle}. Details omitted.
\end{proof}

\end{document}
